% =====================================================================
% VERSION ENSAE, REDUCTION SFE (phase 3, 1er octobre 2026) : enseignements du stage.
% Copie a texte constant de chapitres/19_enseignements_stage.tex.
% Ne changent que les titres et les renvois repares, listes dans
% reduction/phase3.md.
% Lu par main_ensae.tex seul : ne jamais le faire lire par main_institut.tex,
% qui garde chapitres/.
% =====================================================================

% =====================================================================
% ENSEIGNEMENTS DU STAGE
%
% CHAPITRE PROPRE À LA VERSION ENSAE. Seul main_ensae.tex l'appelle : le barème
% de l'école note sur 6 l'analyse et la réflexion sur le contexte et les outils,
% demande de préciser la part d'influence du maître de stage et des tiers, et
% attend un recul sur l'expérience. Un mémoire d'actuariat n'a rien à porter de
% tout cela, et l'Institut ne le demande pas.
%
% RÉÉCRIT EN ENTIER LE 18 SEPTEMBRE 2026, sur consigne de Kélian. Trois règles
% commandent cette version et il ne faut pas les défaire :
%   - registre strictement factuel et concis. Aucune anecdote, aucune mise en
%     scène, aucun récit d'erreur : la version antérieure portait un paragraphe
%     « sept fois j'ai rédigé... et sept fois la sortie l'a démentie » et une
%     série d'anecdotes de vérification de chiffres, tous retirés ;
%   - PAS DE MENTION D'OUTILS D'IA GÉNÉRATIVE DANS CE CHAPITRE. Le paragraphe
%     d'usage qui y existait depuis le 16 septembre en a été retiré le 18, et il
%     n'a pas à y revenir : ce chapitre porte le recul sur la mission, pas la
%     déclaration d'outillage.
%     ATTENTION, CETTE RÈGLE A CHANGÉ DE SENS LE 21 SEPTEMBRE. Elle ne vaut plus
%     que pour CE fichier. Le mémoire, lui, DÉCLARE cet usage, parce que le § 4.5
%     des Recommandations Jury de l'Institut exige une transparence totale et
%     prévoit l'ajournement << sans besoin d'une autre motivation >> en cas de
%     manquement. La déclaration vit donc à deux endroits, et il ne faut les
%     supprimer ni l'un ni l'autre :
%       - le paragraphe d'usage en fin de résumé (chapitres/01_resume.tex),
%         exigé au § 5.1.c ;
%       - l'annexe d'inventaire (chapitres/23_inventaire_ia.tex), exigée au
%         § 5.1.h.
%   - encadrement nommé et traité sobrement : Hugo Rapior en entreprise,
%     Caroline Hillairet à l'ENSAE. Les répondants au formulaire de relecture,
%     eux, ne sont PAS nommés, décision de Kélian du 9 septembre.
%
% La première personne est admise ici, et ici seulement : le barème de l'école
% exige le recul personnel. Le reste du mémoire reste à la voix impersonnelle.
%
% AUCUN RENVOI VERS LES ANNEXES DÉPORTÉES. Les trois \ref vers sec:ann-harnais
% et celui vers sec:ann-intervalle95 ont disparu avec cette réécriture : leurs
% étiquettes n'existent plus depuis le déport des annexes B, D et E.
% =====================================================================

\chapter*{Enseignements du stage}
\phantomsection
\addcontentsline{toc}{chapter}{Enseignements du stage}
\markboth{Enseignements du stage}{}
\label{ens:ch}
\label{part:stage}

% HARNAIS-HORS-SECTION: chapitre de recul sur la mission. Il ne publie aucun
% resultat chiffre : les seules valeurs qu'il evoque sont etablies et sourcees
% dans le corps du memoire, et il n'en reprend aucune en chiffres.


% Reprise le 1er octobre 2026 tard le soir (« une ecriture plus humaine ») : registre cale sur des
% memoires d'actuariat publies avant 2022 (style/registre_humain.md), « je » du recul personnel
% employe plus franchement. Faits inchanges, aucun nombre ajoute.

\section*{Cadre de la mission}
\label{ens:cadre}

J'ai effectué mon stage chez Nexialog Consulting, cabinet de conseil en risques pour le secteur
financier. La mission portait sur la quantification du besoin de capital associé à la
non-conformité au règlement DORA, qui est le sujet de ce mémoire. Quatre contraintes propres à ce
cadre ont orienté le travail.

\subsection*{Livrables}
Outre le mémoire, le stage a donné lieu à un livrable pour le cabinet : l'analyse actuarielle de
l'édition $2026$ d'une étude de marché sur la cyberassurance en France. Je l'ai co-écrite avec mon
maître de stage, et le cabinet l'a publiée \citep{RapiorKaddouri2026}. Cette étude décrit un
marché observé, alors que je construis un modèle sur des états hypothétiques. Elle constate
également, du point de vue de praticiens qui suivent ce marché, que l'exposition cyber ne se traduit
pas en besoin de capital réglementaire, ce qui est la prémisse du mémoire.

\subsection*{Donnée sous licence}
La base de pertes opérationnelles utilisée pour calibrer la sévérité est sous licence et ne peut pas
être redistribuée. Elle n'est donc pas versionnée avec le code, et un lecteur ne peut pas rejouer la
chaîne de calcul depuis la donnée brute. Ce sont les \emph{sorties} des scripts qui sont versionnées,
et le contrôle porte sur la correspondance entre ces sorties et le document. La traçabilité remplace
ainsi la reproductibilité complète, ce que le rapport indique.

\subsection*{Anonymisation des entités}
L'application à quatre assureurs réels est anonymisée. Les sources sont des rapports publics, mais
le modèle n'utilise que le bilan et l'état de conformité y est supposé : nommer une entité
reviendrait à lui attribuer une non-conformité que rien n'observe. La table de correspondance existe
cependant, dans un script, de sorte que le travail reste vérifiable.

\subsection*{Gel de la calibration}
J'ai gelé la calibration à mi-parcours. Il s'agit d'une décision de conduite de projet : un modèle
qui bouge encore, une fois de nombreux résultats écrits, invalide en silence ce qui a déjà été
rédigé. Le gel sépare la correction d'erreur, toujours autorisée, de la recalibration, qui ne l'est
plus. Son prix figure au chapitre~\ref{ch:robustesse}, où des écarts connus sont publiés chiffrés
sans être corrigés.

\section*{Enseignements de l'ENSAE}
\label{ens:ensae}

Le sujet a mobilisé plusieurs enseignements à la fois. La statistique des valeurs extrêmes sert à
modéliser la sévérité, et les modèles à variable latente gaussienne la dépendance entre états de
conformité. Les processus stochastiques décrivent la dynamique de ces états ainsi que le processus
de branchement. La statistique inférentielle fournit les tests d'adéquation et de couverture, les
méthodes de Monte-Carlo les lois de perte simulées, la théorie des jeux coopératifs l'attribution
par canal, et l'assurance non-vie la charge de capital elle-même.

Deux apports ont pesé davantage que les autres. Le premier est l'économétrie de l'identification :
un paramètre mal estimé appelle davantage de données, un paramètre \emph{non identifié} une autre
approche. Reconnaître qu'un paramètre relevait du second cas a changé la méthode
retenue (chapitre~\ref{ch:identifiabilite}). Le second est la lecture de l'indice de queue comme
condition d'existence des moments, qui m'a permis d'écarter deux variantes de modèle au motif qu'un
moment y cesse d'exister. La programmation scientifique a compté autant que ces deux apports,
car maintenir une chaîne de scripts reproductible pendant plusieurs mois relève de l'ingénierie
logicielle.

\section*{Encadrement et tiers}
\label{ens:encadrement}

L'encadrement a été double. En entreprise, Hugo Rapior, mon maître de stage, a tenu avec moi un
point hebdomadaire, préparé sous la forme d'un jeu de diapositives. À l'ENSAE, Caroline Hillairet,
ma tutrice, a suivi le travail lors d'échanges plus espacés.

Leurs demandes n'ont pas porté sur des calculs. Celles de mon maître de stage m'ont conduit à
retirer des lectures que la mesure ne soutenait pas, et à préciser un vocabulaire dont l'imprécision
masquait une borne prise à l'envers. Celles de ma tutrice ont porté sur le traitement de
l'incertitude et ont donné lieu à deux sections du mémoire, le quantile prédictif et la sensibilité
aux probabilités de propagation. J'en retiens qu'une relecture extérieure est surtout utile parce
qu'elle interroge un énoncé hors du contexte qui l'avait justifié.

Le traitement des processus auto-excités appelle une mention particulière, puisque c'est le domaine
de recherche de ma tutrice. L'écartement de cet outil a été testé contre ses propres variantes de
noyau, puis reformulé en cours de travail en équivalence observationnelle suivie d'un choix de
parcimonie (chapitre~\ref{ch:robustesse}).

\subsection*{Sollicitations extérieures}
J'ai lancé deux sollicitations extérieures. J'ai d'abord demandé à un collègue du cabinet un second
codage en aveugle du corpus de rapports post-incident. Ce codage n'était pas reçu à la date de
rédaction, et le dispositif d'accueil, prêt et testé, ne produit aucune donnée en son absence.

J'ai ensuite sollicité trois praticiens du secteur pour une relecture critique des affirmations
qualitatives du modèle, et les trois ont répondu. Le premier n'en a contredit aucune, mais son
exemple admet deux lectures opposées : il illustre la frontière d'identifiabilité sans la lever. Le
deuxième en a contredit deux et a avancé la cause commune comme explication concurrente de la
contagion, tandis que le troisième conteste les deux mêmes affirmations, ce qui répète le désaccord
sans l'élargir. Ces retours sont déclarés sans être intégrés au modèle, ce qui supposerait de
recalibrer, et ils ne déplacent pas les bornes publiées, calculées sans fixer la direction de
propagation. Les répondants ne sont pas nommés, et le détail figure au
\S\ref{sec:relecture-praticien}.

\section*{Bilan personnel}
\label{ens:retiens}

\subsection*{Méthodes quantitatives et objet réglementaire}
La mission a consisté à traduire un texte réglementaire, qui énonce des exigences d'organisation, en
paramètres d'une loi de perte. Cela suppose de décider ce que le texte permet d'affirmer et ce que
la donnée permet d'estimer, puis de maintenir la distinction entre les deux dans chaque énoncé
publié.

\subsection*{Paramètres estimés, posés et non identifiés}
La modélisation de la conformité DORA mêle trois statuts de paramètre. Certains sont estimés sur
données, d'autres sont posés faute de source, et la direction de propagation entre domaines de
contrôle n'est identifiée par aucune des sources disponibles. Afficher chaque paramètre sous son
statut permet à un tiers de discuter le résultat.

\subsection*{Vérifiabilité des résultats}
En conseil, un résultat doit pouvoir être contrôlé par quelqu'un qui n'a pas fait le calcul.
Rattacher chaque nombre publié au calcul qui le produit répond à cette exigence, et ce travail a
occupé une part importante de mon stage. C'est la différence la plus nette que j'aie observée entre
un exercice académique et une mission.

\section*{Recul sur la démarche}
\label{ens:autrement}

Avec le recul, j'aurais mis en place le dispositif de vérification dès le début du stage : construit
à mi-parcours, il m'a obligé à rattacher après coup chaque section du document à ses scripts.
J'aurais aussi gelé la calibration plus tôt, car plusieurs semaines de recalibrations partielles ont
fait bouger des nombres déjà écrits avant que la décision soit prise.
